\documentclass[11pt]{article}
\usepackage[T1]{fontenc}
\usepackage{lmodern}
\usepackage[margin=1in]{geometry}
\usepackage{amsmath,amssymb,amsthm,mathtools}
\usepackage{microtype,booktabs,enumitem,xurl,xcolor}
\usepackage[colorlinks=true,linkcolor=blue!45!black,citecolor=blue!45!black,urlcolor=blue!45!black]{hyperref}
\newtheorem{theorem}{Theorem}[section]
\newtheorem{lemma}[theorem]{Lemma}
\newtheorem{proposition}[theorem]{Proposition}
\newtheorem{corollary}[theorem]{Corollary}
\theoremstyle{remark}\newtheorem{remark}[theorem]{Remark}
\numberwithin{equation}{section}
\newcommand{\R}{\mathbb R}\newcommand{\C}{\mathbb C}
\newcommand{\E}{\mathbb E}\newcommand{\Prob}{\mathbb P}
\newcommand{\HH}{\mathcal H}\newcommand{\op}{\mathrm{op}}
\DeclareMathOperator{\Tr}{Tr}\DeclareMathOperator{\TV}{TV}
\DeclareMathOperator{\sgn}{sign}\DeclareMathOperator{\diam}{diam}
\newcommand{\ip}[2]{\langle #1,#2\rangle}
\newcommand{\norm}[1]{\lVert #1\rVert}
\allowdisplaybreaks[1]
\setlist{itemsep=3pt,topsep=5pt}
\hypersetup{pdftitle={Spectral chart selection and average trace-class distortion},pdfauthor={Myeongjin Kang}}
\title{Spectral chart selection and\\average trace-class distortion}
\author{Myeongjin Kang}
\date{October 18, 2026}
\begin{document}
\maketitle
\begin{abstract}
We construct a sequence of genuine finite negative-type metrics with $n$
points whose first uniform-average distortion into the trace class
$S_1(H)$ is at least $c\sqrt{\log n}/(\log\log n)^{9/2}$, uniformly over
all complex Hilbert spaces $H$. The construction uses the multichart
Hilbert kernel of the pinned uniform-sparsest-cut source. Its scalar
coordinate selection is replaced by a dimension-independent matrix
selection: universal scalar moment matching gives spectral cutoffs for
all matrix tuples, and a positive Gaussian dual pairing averages these
cutoffs without requiring them to commute. We track the enlarged chart
number through the kernel, rounding mesh and triangle repair. Finite-rank
compression, Hermitian dilation and uniform copies preserve the full
statement. The result is new to this project; a bounded search found no
antecedent, but literature priority is not claimed.
\end{abstract}

\section{Statement, conventions and source dependence}
For a complex Hilbert space $H$, write $S_1(H)$ for its trace-class
operators, with norm $\norm{A}_1=\Tr|A|$. A finite metric $d$ has
\emph{negative type} if there are real Hilbert vectors $z_x$ with
$d(x,y)=|z_x-z_y|^2$. We require $d(x,y)>0$ for $x\ne y$ and every triangle
inequality; a squared Hilbert distance without these properties is not
called a metric here. All pair averages are over two independent uniform
points, or equivalently over all ordered pairs including the diagonal.

\begin{theorem}\label{thm:main}
There are an absolute constant $c>0$ and genuine finite negative-type
metrics $(V_j,d_j)$, with $n_j=|V_j|\to\infty$, such that for every complex
Hilbert space $H$ and every map $F:V_j\to S_1(H)$ satisfying
$\norm{F(x)-F(y)}_1\le d_j(x,y)$ for all $x,y$, one has
\begin{equation}\label{eq:main}
 \frac1{n_j^2}\sum_{x,y\in V_j}\norm{F(x)-F(y)}_1
 \le c^{-1}\frac{(\log\log n_j)^{9/2}}{\sqrt{\log n_j}}
       \frac1{n_j^2}\sum_{x,y\in V_j}d_j(x,y).
\end{equation}
There is no dimension, separability or rank restriction on $H$ or $F$.
\end{theorem}
This also gives the displayed lower bound on bi-Lipschitz distortion
within negative-type metrics. For general finite metrics it is weaker
than the classical cotype grid bound recalled in
Section~\ref{sec:comparisons}; no new general-metric bi-Lipschitz bound
is claimed. It is an existence result, with no efficient or rational construction
claim. We write $C$ for absolute constants, including $c^{-1}$ after
adjusting $c$, and all logarithms are natural.

Our geometric input is OpenAI's \emph{Near-square-root logarithmic
integrality gaps for uniform sparsest cut} \cite{source}, at the immutable
revision specified in Appendix~\ref{sec:source}. Its theorems are taken as
inputs; this paper does not audit their truth. We use its explicit common
kernel, chart extension, angle repair, interface coarea and copy mechanism.
The original scalar contraction cannot simply be asserted for trace
norms: its good charts are selected separately for scalar coordinates.
Section~\ref{sec:selector} supplies the replacement. The other operator
and probability tools are standard, and are proved here in the particular
forms needed.

The chart number and exponent were chosen by the method before the theorem
scope was frozen; the freeze preceded the complete proof, not the method.
Specifically, set, for a sufficiently large integer $m$,
\begin{equation}\label{eq:parameters}
\begin{gathered}
 N=m^6,\quad q=\lceil\log N\rceil,\quad S=m^{10q},\qquad
 a=m^{-40},\quad b=m^{40},\\
 \varepsilon=m^{-200}/S,\quad r=m^{-8},\quad
 \tau=m^{-2000}/S^2,\quad \lambda=m^{-2}.
\end{gathered}
\end{equation}
The denominator exponent $9/2$ is method-shaped, not an asserted optimum.
Restricting \eqref{eq:main} to the diagonal copy of $\ell_1$ in $S_1$ gives
a weaker commutative bound than the source's exponent $3$; our theorem
does not supersede that result. Section~\ref{sec:comparisons} treats the
closest alternative operations, including the sharp metric-cotype grid
bound and the pair-average inequality implicit in the cube-quotient
Talagrand argument. We make no claim about the trace-class embedding of
Heisenberg groups or the general quantum-expansion metric question.

\section{A variance-normalized operator moment bound}\label{sec:wick}
All matrices in Sections~\ref{sec:wick}--\ref{sec:selector} are finite
complex matrices. The notation $A\le B$ means Loewner order.
\begin{lemma}\label{lem:wick}
If $X_1,\ldots,X_m$ are selfadjoint and $\sum_jX_j^2\le I$, then for
independent standard real Gaussian variables $g_j$ and every integer
$q\ge1$,
\begin{equation}\label{eq:wick}
 0\le \E\Big(\sum_jg_jX_j\Big)^{2q}\le (2q-1)!!\,I.
\end{equation}
\end{lemma}
\begin{proof}
Isserlis' identity, including repeated indices, expresses the expectation
exactly as a sum over all pairings $\pi$ of the $2q$ positions. For a fixed
pairing, let $W_\pi$ be the sum of words obtained by assigning a free index
in $[m]$ to each pair and putting its $X$ at both positions.

On a Hilbert space $K$ carrying the matrices, define
$V:K\to\C^m\otimes K$ by $Vv=\sum_i e_i\otimes X_iv$. Since
$V^*V=\sum_iX_i^2\le I$, both $V$ and $V^*$ are contractions.
Process the positions of the word from right to left. At the first
encounter with a pair, apply $V$ to the $K$ factor and keep the new
$\C^m$ factor as its index register. At its second encounter, permute
that register next to $K$ and apply $V^*$. Tensor each operation with the
identity on all other open registers. Permutations are unitary, so every
operation is contractive. This applies to crossing pairings as well as
nested ones. At the end no registers remain.

Expanding the composition shows that opening a register sums its index,
closing it enforces equality with the matched occurrence, and the order
of the $X$ factors is the original word order. Thus the resulting map is
$W_\pi$ and $\norm{W_\pi}_{\op}\le1$. No positivity of an individual
$W_\pi$ is used. There are $(2q-1)!!$ pairings. Their sum has norm at most
that number, and the sum is positive because it is the expectation of
an even power of a selfadjoint matrix. This proves \eqref{eq:wick}.
\end{proof}
Operator-valued Khintchine and Wick estimates belong to the classical
noncommutative probability literature; see, for example,
Haagerup--Pisier \cite{HP} and Buchholz \cite{Buchholz}. We credit this
phenomenon as a standard tool. The proof above supplies the exact Loewner
form used here rather than inferring it from a trace-moment estimate.

\section{Directions universal over all matrix dimensions}\label{sec:design}
\begin{lemma}\label{lem:design}
For the parameters in \eqref{eq:parameters}, there are vectors
$g_{s,i}\in\R^m$, $1\le s\le S$, $1\le i\le N$, such that
$u_{s,i}=g_{s,i}/\sqrt m$ satisfy
\begin{equation}\label{eq:moment}
 \frac1{SN}\sum_{s,i}\Big(\sum_j u_{s,ij}X_j\Big)^{2q}
 \le (4q/m)^q I
\end{equation}
for every selfadjoint tuple of every finite dimension with
$\sum_jX_j^2\le I$. They also satisfy, for every chart $s$,
\begin{gather}
 \sqrt m/2\le |g_{s,i}|\le2\sqrt m,\qquad
 \tfrac12 I\le N^{-1}\sum_i g_{s,i}g_{s,i}^{T}\le2I,\label{eq:dirnorm}\\
 \sup_{|w|\le m^{42}}\left|
   N^{-1}\sum_i g_{s,i}\sin\ip{g_{s,i}}w-we^{-|w|^2/2}\right|
 \le 2/\sqrt m.\label{eq:fourier}
\end{gather}
The directions within each chart are pairwise nonparallel.
\end{lemma}
\begin{proof}
Sample all $SN=m^{10q+6}$ vectors independently from standard Gaussian
measure. For an ordered word $\alpha=(j_1,\ldots,j_{2q})$, truncate
$G^\alpha=G_{j_1}\cdots G_{j_{2q}}$ on $|G|>2\sqrt m$. Its absolute
value is at most $B=(2\sqrt m)^{2q}$. Gaussian radial tails give
\[
 \big|\E G^\alpha-\E[G^\alpha\mathbf1_{|G|\le2\sqrt m}]\big|
 \le(\E|G|^{4q})^{1/2}\Prob(|G|>2\sqrt m)^{1/2}
 \le (Cm)^q e^{-cm}<m^{-3q},
\]
for all sufficiently large $m$, since $q=O(\log m)$.
Hoeffding's bounded-variable inequality at tolerance $m^{-3q}\le B$
shows that the empirical truncated moment differs from its expectation
by more than $m^{-3q}$ with probability at most
\[
 2\exp\left(-c\frac{m^{10q+6}m^{-6q}}{16^q m^{2q}}\right)
 =2\exp(-c m^{2q+6}/16^q).
\]
A union bound over $m^{2q}$ ordered words still tends to zero.
Also $SN\Prob(|G|\notin[\sqrt m/2,2\sqrt m])\to0$, because
$\log(SN)=O((\log m)^2)$. On the latter event the sampled truncated
moments equal the untruncated ones. Hence all empirical moments differ
from their exact Gaussian values by at most $2m^{-3q}$ simultaneously.

For any admissible tuple, $X_j^2\le\sum_kX_k^2\le I$, so
$\norm{X_j}_{\op}\le1$ and the norm of each word $X_\alpha$ is at most
one. The factor $m^{-q}$ from $u=g/\sqrt m$ therefore makes the
empirical-minus-Gaussian operator moment error at most
\[
 m^{-q}\,m^{2q}\,2m^{-3q}=2m^{-2q}.
\]
The error is selfadjoint. Lemma~\ref{lem:wick} bounds the exact normalized
Gaussian moment by $(2q-1)!!m^{-q}I$, and
$(2q-1)!!m^{-q}+2m^{-2q}\le(4q/m)^q$ for large $m$. This proves
\eqref{eq:moment}. The event fixed only scalar coefficients before a
matrix tuple or its dimension was chosen; this is the source of the
universal quantifier.

For clarity, the additional geometric direction events are the first two
parts of the source directions lemma \cite{source}. Its probabilistic
proof bounds the Fourier/covariance failures by
\[
 \exp[-cm^4+O(m\log m)+\log S]
\]
and the norm failures by $SN e^{-cm}$. These still tend to zero for our
$\log S=O((\log m)^2)$. The Fourier net has radius $m^{-3}$ on a ball of
radius $m^{42}$; a $1/8$ sphere net controls vector and covariance norms.
The truncated empirical Fourier function is $4m$-Lipschitz and its exact
Gaussian counterpart is $2$-Lipschitz, which gives the displayed uniform
error after net interpolation. Nonparallelism is almost sure. Intersect
these events with the moment event and fix one realization. The source's
scalar majority-good-chart assertion is not used.
\end{proof}
Matching scalar polynomial moments is a standard design principle; exact
averaging sets and spherical designs provide classical context
\cite{SZ}. We need only the elementary approximate Gaussian matching
just proved, not a new design existence theorem or a sphere-design bound.

Set $(J_s\theta)_i=\ip{u_{s,i}}\theta$ and
$J_s^+=(J_s^TJ_s)^{-1}J_s^T$. The covariance conclusion gives
\begin{equation}\label{eq:pseudoinverse}
 J_s^+J_s=I,\quad
 \norm{J_s^+}_{\op}\le2\sqrt{m/N},\quad
 |J_s^+e_i|\le4m/N.
\end{equation}
Here the last estimate uses $|u_{s,i}|\le2$ and
$\norm{(J_s^TJ_s)^{-1}}_{\op}\le2m/N$.

\section{Spectral chart selection}\label{sec:selector}
We reserve $M_s$ for the positive matrix controlling a chart and $P_s$
for its spectral selector. Geometric chart complements will be $R_s$,
and their extended Hilbert features will be $\mathcal P_s$.

\begin{lemma}\label{lem:selector}
Suppose the directions satisfy \eqref{eq:moment}. Let $A_j,A_{s,j}$ and
$b_{s,i}$ be selfadjoint matrices such that
\[
 A_{s,j}=\sum_i u_{s,ij}b_{s,i},\qquad
 \norm{b_{s,i}}_1\le t,\qquad
 \sum_j\norm{A_{s,j}-A_j}_1\le e
\]
for every $s$, where $t,e\ge0$. Then
\begin{equation}\label{eq:selector}
 \E_g\norm{\sum_jg_jA_j}_1
 \le 2\sqrt{8\exp(1)}\,Nt\sqrt{q/m}+2e.
\end{equation}
The bound is independent of the matrix dimension.
\end{lemma}
\begin{proof}
Put $A(g)=\sum_jg_jA_j$, $U(g)=\sgn(A(g))$ with $\sgn(0)=0$, and
$X_j=\E[g_jU(g)]$. These are selfadjoint. For each vector $v$, Bessel's
inequality in $L_2(\gamma_m;\C^D)$ gives
\[
 \sum_j|X_jv|^2\le\E|U(g)v|^2\le|v|^2.
\]
Thus $\sum_jX_j^2\le I$. The essential positive pairing is
\begin{equation}\label{eq:positive}
 Y:=\sum_jA_jX_j=\E[A(g)\sgn(A(g))]=\E|A(g)|\ge0,
 \qquad \Tr Y=\E\norm{A(g)}_1.
\end{equation}
No positivity is asserted for the individual products $A_jX_j$.

Set
\[
 Z_{s,i}=\sum_j u_{s,ij}X_j,\qquad
 M_s=\Big(\sum_iZ_{s,i}^{2q}\Big)^{1/q}.
\]
Operator concavity of $T\mapsto T^{1/q}$ and \eqref{eq:moment} imply
\[
 S^{-1}\sum_s M_s
 \le\Big(S^{-1}\sum_{s,i}Z_{s,i}^{2q}\Big)^{1/q}
 \le N^{1/q}(4q/m)I\le4\exp(1)(q/m)I.
\]
Let $h=8\exp(1)q/m$ and $P_s=\mathbf1_{[0,h]}(M_s)$. Spectral calculus
on each $M_s$ gives $I-P_s\le M_s/h$, whence
\begin{equation}\label{eq:coverage}
 \sum_sP_s\ge(S/2)I.
\end{equation}
Also $Z_{s,i}^{2q}\le M_s^q$. Operator monotonicity of $T^{1/q}$ gives
$Z_{s,i}^2\le M_s$, and hence
\begin{equation}\label{eq:compression}
 \norm{Z_{s,i}P_s}_{\op}^2
 =\norm{P_sZ_{s,i}^2P_s}_{\op}\le h.
\end{equation}
For completeness, both fractional-power properties used here follow, for
$0<\alpha<1$, from
\[
 T^\alpha=\frac{\sin(\pi\alpha)}\pi
     \int_0^\infty T(T+sI)^{-1}s^{\alpha-1}\,ds.
\]
The inverse reverses Loewner order and is operator convex. Therefore
$I-s(T+sI)^{-1}$ is operator monotone and concave; integration preserves
these properties. Continuity handles zero eigenvalues, and $\alpha=1$
is immediate. Different $M_s$ are never diagonalized simultaneously.

Since $Y\ge0$, \eqref{eq:coverage} can be paired with $Y$ under the trace:
\begin{align*}
 (S/2)\Tr Y
 &\le\sum_s\Tr(YP_s)\\
 &=\sum_{s,i}\Tr(b_{s,i}Z_{s,i}P_s)
   +\sum_{s,j}\Tr((A_j-A_{s,j})X_jP_s)\\
 &\le SNt\sqrt h+Se.
\end{align*}
Here $\Tr(YP_s)=\Tr(Y^{1/2}P_sY^{1/2})$ is real and nonnegative;
the expanded summands can be nonreal, so the final line bounds their
absolute values by trace duality. We use \eqref{eq:compression} and
$\norm{X_jP_s}_{\op}\le1$. Division by $S/2$ and \eqref{eq:positive}
prove \eqref{eq:selector}.
\end{proof}

\section{The geometric construction with the larger chart family}
\label{sec:geometry}
This section uses the explicit operations in the pinned source
\cite{source}, with every changed parameter tracked. We do not invoke its
fixed-$S=m^3$ metric proposition unchanged.

\subsection{Common kernel and exact chart extensions}
Let $E(t)=(\cos t,\sin t)$ and let $\gamma_m$ be Gaussian measure. Use the
real Hilbert sum
\[
 \HH=L_2([a,b]\times\R^m,d\sigma\,d\gamma_m;\R^2)
 \oplus\bigoplus_{s,i}L_2([a,b],d\sigma/\sigma;\R^2).
\]
Define $p(\theta)$ to have first component
$m^{1/4}E(\ip g\theta/\sigma)$ and component $(s,i)$ equal to
$\sqrt\varepsilon E(\ip{g_{s,i}}\theta/\sigma)$. These are exactly the
source common-feature formulas, with our $S$ and $\varepsilon$. They give
\begin{align}
 \ip{p(\theta)}{p(\eta)}
 &=\sqrt m\int_a^b e^{-|\eta-\theta|^2/(2\sigma^2)}\,d\sigma\notag\\
 &\quad+\varepsilon\sum_{s,i}\int_a^b
       \cos(\ip{g_{s,i}}{\eta-\theta}/\sigma)\,\frac{d\sigma}{\sigma}.
       \label{eq:kernel}
\end{align}
For $|\theta|,|\eta|\le3\sqrt m$, put
$c_*=2\int_0^\infty(1-e^{-1/(2t^2)})\,dt>0$. Scaling the Gaussian
integral and bounding the omitted ends show that
\begin{equation}\label{eq:distanceerror}
 \big||p(\theta)-p(\eta)|^2-c_*\sqrt m|\theta-\eta|\big|
 \le 2\sqrt m a+36m^{3/2}/b+4\varepsilon SN\log(b/a)
 \le m^{-37}.
\end{equation}
In particular, $\varepsilon SN=m^{-194}$ despite the larger chart family.
For $j=0,1,2$, the squared derivative bounds from the Gaussian and
auxiliary blocks are, respectively,
\[
 \sqrt m\,b\,a^{-2j}\E|G|^{2j},\qquad
 \varepsilon SN\log(b/a)(2\sqrt m/a)^{2j}.
\]
Thus $|p|,\norm{Dp}_{\op},\norm{D^2p}_{\op}\le m^{110}$, with no
$\sqrt S$ loss.

Let $B_s^0(\theta)e_i$ be supported in block $(s,i)$ and equal there to
\[
 (m/N)\varepsilon^{-1/2}E'(\ip{g_{s,i}}\theta/\sigma).
\]
Define
\begin{equation}\label{eq:B}
 B_s=B_s^0(I-J_sJ_s^+)+Dp\,J_s^+,
 \qquad R_s=I-J_sJ_s^+.
\end{equation}
Then $B_sJ_s=Dp$ exactly. This is the source chart-derivative construction,
renamed to avoid collision with the matrices of Lemma~\ref{lem:selector}.
Its profile remains
\begin{equation}\label{eq:profile}
 |\ip{B_s(\theta)e_i}{p(\eta)}|\le C(m/N)\log m
 \quad (|\theta|,|\eta|\le3\sqrt m).
\end{equation}
Here is the explicit parameter check. In the main profile the factors
$\sqrt\varepsilon$ and $\varepsilon^{-1/2}$ cancel, leaving
$(m/N)\int_a^b\sin(\ip{g_{s,i}}{\eta-\theta}/\sigma)d\sigma/\sigma$.
The residual $Dp-B_s^0J_s$ pairs with $p(\eta)$ through the Fourier
error \eqref{eq:fourier}, because $|\eta-\theta|/\sigma\le6\sqrt m/a<m^{42}$.
Its main part is at most $2\log(b/a)$, and its auxiliary gradient is at most
$2\varepsilon SN\sqrt m/a\le1$. Multiplying the residual by $J_s^+e_i$
and using \eqref{eq:pseudoinverse} proves \eqref{eq:profile}.

The columns of $B_s^0$ have disjoint Hilbert supports. Even the looser
bound by $\sqrt N$ times the largest column norm gives
\[
 \norm{B_s^0}_{\op}\le\sqrt N(m/N)\varepsilon^{-1/2}\sqrt{\log(b/a)},
 \quad
 \norm{DB_s^0}_{\op}\le\sqrt N(m/N)\varepsilon^{-1/2}
          \sqrt{\log(b/a)}\,2\sqrt m/a.
\]
Both acquire only a factor $\sqrt S$ relative to the original powers of
$m$. Moreover
\[
 DB_s=DB_s^0(I-J_sJ_s^+)+D^2p\,J_s^+.
\]
It follows that $\norm{B_s}_{\op},\norm{DB_s}_{\op}\le\sqrt S m^{300}$.
Thus a common bound for all five strong norms is $\sqrt S m^{300}$,
while the first three still satisfy the sharper $m^{110}$ bound.

\subsection{Rounding, angles and repair}
Let $U=(-2,2)^m$, and away from the threshold hyperplanes define
\[
 x_s(\theta)_i=\tau\left\lfloor\ip{u_{s,i}}\theta/\tau\right\rfloor,
 \quad v_s(\theta)=(s,x_s(\theta)),\quad
 V=\{v_s(\theta):\theta\in U\text{ off those hyperplanes}\}.
\]
Set
\[
 \mathcal P_s(x)=p(J_s^+x)+B_s(J_s^+x)R_sx,
 \qquad \mathcal P_v=\mathcal P_s(x)\ (v=(s,x)).
\]
Rounding gives
\begin{equation}\label{eq:round}
 |x-J_s\theta|\le\sqrt N\tau,\quad
 |J_s^+x-\theta|\le2\sqrt m\tau,\quad
 |R_sx|\le\sqrt N\tau.
\end{equation}
Consequently, the feature error satisfies
\[
 |\mathcal P_{v_s(\theta)}-p(\theta)|\le\sqrt S m^{305}\tau.
\]
The resulting squared-distance error is at most
$CS^{1/2}m^{415}\tau\le S m^{610}\tau$. By \eqref{eq:parameters},
$S m^{610}\tau=m^{-1390}/S\ll r$, so
\begin{equation}\label{eq:rounded}
 \big||\mathcal P_{v_s(\theta)}-\mathcal P_{v_t(\eta)}|^2
       -c_*\sqrt m|\theta-\eta|\big|\le r/10.
\end{equation}
The three-distance polarization identity and the Euclidean triangle
inequality then give every angle a lower bound $-r$.

For same-chart labels $x,x'$, every point $\xi$ on their segment satisfies
the last two bounds of \eqref{eq:round}, with $|J_s^+\xi|\le3\sqrt m$.
Differentiation gives the exact identity
\begin{equation}\label{eq:DP}
 D\mathcal P_s(\xi)v
 =B_s(J_s^+\xi)v+DB_s(J_s^+\xi)[J_s^+v]R_s\xi.
\end{equation}
The second term has operator norm at most $2\sqrt S m^{300.5}\tau$.
When paired with $\mathcal P_y$, it contributes at most
$C\sqrt S m^{411}\tau$. Replacing $\mathcal P_y$ by its common feature
in the first term costs at most $CSm^{605}\tau$. These are distinct error
sources; both are bounded by $CSm^{610}\tau\ll(m/N)\log m$.
Equation~\eqref{eq:profile}, applied also to $\mathcal P_z$, and integration
along the segment therefore give, with $L=C_0\log m$,
\begin{equation}\label{eq:angle}
 |\ip{\mathcal P_v-\mathcal P_z}{\mathcal P_y-\mathcal P_z}|
 \le L(m/N)\norm{x-x'}_1
 \quad(v=(s,x),\ z=(s,x')).
\end{equation}

Set $\delta(v,w)=L(m/N)\norm{x-x'}_1$ within a chart and $\delta=\infty$
between charts. The source triangle-repair lemma applies to any finite
chart partition: if the angles are at least
$-\min\{r,\delta(v,z),\delta(y,z)\}$, it appends Hilbert vectors $q_v$
with
\begin{equation}\label{eq:repair}
 |q_v-q_w|^2=8\int_0^r(1-e^{-\delta(v,w)/\ell})\,d\ell\le8r
\end{equation}
such that $d(v,w)=|\mathcal P_v-\mathcal P_w|^2+|q_v-q_w|^2$ obeys
all triangle inequalities. The global bound and \eqref{eq:angle} give
exactly its hypothesis. Its positive kernel is
$4\int_0^r e^{-\delta(v,w)/\ell}d\ell$, block diagonal across charts;
this is the credited Schoenberg-kernel repair, not a new operation.
Distinct same-chart labels have $\delta>0$, and between charts
\eqref{eq:repair} equals $8r$. Thus $d$ is already a genuine negative-type
metric on the label set.

Taking $y=v$ in \eqref{eq:angle} bounds the original squared feature
distance by $\delta$. For $0<\delta\le r$, the repair cost is at most
$8\delta(1+\log(r/\delta))$. Therefore adjacent labels satisfy
\begin{equation}\label{eq:adjacent}
 \norm{x-x'}_1=\tau\quad\Longrightarrow\quad
 d((s,x),(s,x'))\le C\tau(m/N)(\log m)^3,
\end{equation}
because
$\log(r/(L(m/N)\tau))=1997\log m+2\log S-\log(C_0\log m)
=O((\log m)^2)$. The macro bound is
\begin{equation}\label{eq:macro}
 |d(v_s(\theta),v_t(\eta))-c_*\sqrt m|\theta-\eta||\le Cr.
\end{equation}

At most $H_0=N(1+8\sqrt m/\tau)$ threshold hyperplanes meet $U$ in one
chart. The arrangement bound, proved by induction on the number of
hyperplanes, gives at most $\sum_{j=0}^m\binom{H_0}j\le(H_0+1)^m$
regions. Each label is constant on a region, so
\begin{equation}\label{eq:size}
 |V|\le S(H_0+1)^m\le\exp(Cm(\log m)^2),\qquad \diam(V,d)\le Cm.
\end{equation}
Finally let $T=(-1,1)^m$. For independent uniform $\theta,\eta\in T$,
$\E|\theta-\eta|^2=2m/3$ and $|\theta-\eta|\le2\sqrt m$, whence
$\E|\theta-\eta|\ge\sqrt m/3$. It follows that
\begin{equation}\label{eq:macroaverage}
 \E_{\theta,\eta\in T}d(v_1(\theta),v_1(\eta))\ge cm.
\end{equation}

\section{Trace-norm contraction on the label set}\label{sec:contraction}
First let $F:V\to M_D(\C)$ have selfadjoint values and be nonexpanding
for $d$. Set $f_s(\theta)=F(v_s(\theta))$ almost everywhere on $U$ and
extend by zero outside $U$, assigning zero also on threshold hyperplanes.
Choose a fixed even nonnegative $\rho\in C_c^\infty((-1,1))$ of integral
one and let
\[
 \kappa_\lambda(z)=\prod_{j=1}^m\lambda^{-1}\rho(z_j/\lambda),
 \qquad H_s=f_s*\kappa_\lambda.
\]
These are smooth finite-dimensional matrix-valued functions. For
$\theta\in\overline T$, the convolution kernel is supported inside $U$.
The macro bound \eqref{eq:macro} implies
\begin{align}
 \norm{H_s(\theta)-f_s(\theta)}_1&\le C(m\lambda+r)
                  &&\text{for almost every }\theta\in T,\label{eq:smoothing}\\
 \norm{\partial_jH_s(\theta)-\partial_jH_1(\theta)}_1&\le Cr/\lambda
                  &&(\theta\in T).\label{eq:mismatch}
\end{align}
Indeed, kernel displacements have Euclidean norm at most $\sqrt m\lambda$,
and $\norm{f_s-f_1}_1\le Cr$ pointwise inside $U$. Differentiating the
kernel uses $\int|\partial_j\kappa_\lambda|=\lambda^{-1}\norm{\rho'}_1$.
These are norm inequalities and require no scalar-coordinate decomposition.

\subsection{The pointwise interface coefficient bound}
Fix $s$ and write $u_i=u_{s,i}$. Across the interface
$Q_{i,k}=\{x:\ip{u_i}x=k\tau\}$, let $\Delta_{i,k}(x)$ be the jump
of $f_s$ toward increasing $\ip{u_i}x$. It exists almost everywhere on
the facet. Other interfaces intersect it in codimension at least two,
by nonparallelism. Adjacent labels differ by $\tau e_i$, so
\[
 \norm{\Delta_{i,k}(x)}_1\le C\tau(m/N)(\log m)^3.
\]
Integration by parts on the polyhedral cells gives, entrywise, the
distributional gradient inside $U$. After convolution it is the classical
identity, at every $\theta\in T$,
\begin{equation}\label{eq:coefficients}
 \partial_j H_s(\theta)=\sum_i u_{s,ij}b_{s,i}(\theta),\qquad
 b_{s,i}(\theta)=\frac1{|u_i|}\sum_k\int_{Q_{i,k}\cap U}
 \kappa_\lambda(\theta-x)\Delta_{i,k}(x)\,d\mathcal H^{m-1}(x).
\end{equation}
The cell interiors contribute zero; paired facet terms give the signed
jump and its unit normal. The testing kernel vanishes near $\partial U$,
so there is no outer boundary term.

Let $h_i$ be the density of $\ip{u_i}Z$ when $Z$ has density
$\kappa_\lambda(\theta-\cdot)$. Coarea identifies
\[
 h_i(t)=|u_i|^{-1}\int_{\ip{u_i}x=t}
                  \kappa_\lambda(\theta-x)\,d\mathcal H^{m-1}(x).
\]
As $|u_i|\ge1/2$, some coordinate has modulus at least $1/(2\sqrt m)$.
Conditioning on all other coordinates shows that $h_i$ is a mixture of
translated or reflected one-dimensional kernels at scale at least
$\lambda/(2\sqrt m)$. Hence $\int|h_i'|\le C\sqrt m/\lambda$.
Comparing values and integrals on each interval of length $\tau$ gives
\[
 \tau\sum_k h_i(k\tau)\le1+\tau\int|h_i'|
 \le1+C\tau\sqrt m/\lambda\le2.
\]
The nonnegative kernel in \eqref{eq:coefficients} now proves the
pointwise, rather than merely integrated, bound
\begin{equation}\label{eq:budget}
 \norm{b_{s,i}(\theta)}_1\le C(m/N)(\log m)^3.
\end{equation}
This is the source coarea calculation in trace norm; the matrices
$b_{s,i}$ need not be positive.

Apply Lemma~\ref{lem:selector} with $A_j=\partial_jH_1(\theta)$,
$A_{s,j}=\partial_jH_s(\theta)$, $t=C(m/N)(\log m)^3$ and
$e=Cmr/\lambda=Cm^{-5}$. Since $q=O(\log m)$, it yields
\begin{equation}\label{eq:GaussianGrad}
 \E_g\norm{\sum_jg_j\partial_jH_1(\theta)}_1
 \le C\sqrt m(\log m)^{7/2}\qquad(\theta\in T).
\end{equation}

\subsection{A Banach-valued first-moment cube inequality}
\begin{lemma}\label{lem:cube}
For a finite-dimensional normed space $B$ and a $C^1$ map $H$ on a
neighborhood of $\overline T$,
\[
 \E_{\theta,\eta\in T}\norm{H(\theta)-H(\eta)}_B
 \le\sqrt{\pi/2}\,
       \E_{\theta\in T}\E_g\norm{\sum_jg_j\partial_jH(\theta)}_B.
\]
\end{lemma}
\begin{proof}
Let $G,G'$ be independent standard Gaussian vectors and, for
$0\le t\le\pi/2$, put $G(t)=\cos t\,G+\sin t\,G'$ and
$W(t)=-\sin t\,G+\cos t\,G'$. At fixed $t$, these are independent
standard Gaussian vectors. With $\Phi,\phi$ the Gaussian distribution
function and density, set $\theta_j(t)=2\Phi(G_j(t))-1$. The endpoints
are independent uniform points of $T$, and the derivative is
\[
 \frac d{dt}H(\theta(t))
 =\sum_j2\phi(G_j(t))W_j(t)\partial_jH(\theta(t)).
\]
The normalized coefficients $2\phi(G_j(t))/\sqrt{2/\pi}$ belong to
$[0,1]$. For any vectors $v_j$ and $0\le a_j\le1$, adding independent
noise $\sqrt{1-a_j^2}\,g_j'$ and applying Jensen proves
$\E\norm{\sum a_jg_jv_j}\le\E\norm{\sum g_jv_j}$.
Condition on $G(t)$, use this contraction inequality, and integrate the
norm of the derivative over $[0,\pi/2]$. Since $\theta(t)$ is uniform
for every $t$, the constant is $(\pi/2)\sqrt{2/\pi}=\sqrt{\pi/2}$.
The derivatives are bounded on the compact cube and Gaussian first
moments are finite, so the integrations are justified. This is the
standard Gaussian rotation method \cite{Pisier}, in the norm-valued form
needed here.
\end{proof}
Combining Lemma~\ref{lem:cube}, \eqref{eq:GaussianGrad} and
\eqref{eq:smoothing}, and using $m\lambda=m^{-1}$, proves
\begin{equation}\label{eq:labelcontraction}
 \E_{\theta,\eta\in T}
 \norm{F(v_1(\theta))-F(v_1(\eta))}_1
 \le C\sqrt m(\log m)^{7/2}.
\end{equation}
For arbitrary complex matrix values use the real-linear Hermitian dilation
\[
 \mathcal D(B)=\begin{pmatrix}0&B\\ B^*&0\end{pmatrix},
 \qquad \norm{\mathcal D(B)}_1=2\norm B_1.
\]
The map $\mathcal D(F)/2$ is selfadjoint and preserves every trace-norm
pair distance. Thus \eqref{eq:labelcontraction} holds without a
selfadjointness assumption.

\section{Uniform copies, arbitrary Hilbert spaces and the exponent}
\label{sec:completion}
Let $w_v=\Prob(v_1(\theta)=v)$ for uniform $\theta\in T$, retaining all
vertices including those of weight zero. Set
\[
 M=m^m+m^2|V|,\qquad k_v=1+\lfloor Mw_v\rfloor,\qquad
 n=\sum_vk_v,\qquad \mu_v=k_v/n.
\]
Writing $k_v=Mw_v+e_v$ with $0<e_v\le1$, we obtain
$M\le n\le M+|V|$ and
\[
 \TV(\mu,w)=\tfrac12\sum_v|\mu_v-w_v|
 \le |V|/M\le m^{-2}.
\]
The product laws differ in total variation by at most $2m^{-2}$.
Both the metric and every nonexpanding image distance are at most $Cm$.
Consequently \eqref{eq:labelcontraction} and \eqref{eq:macroaverage}
imply, for every finite-matrix-valued nonexpanding $F$,
\begin{equation}\label{eq:mu}
 \E_{v,w\sim\mu}\norm{F(v)-F(w)}_1\le C\sqrt m(\log m)^{7/2},
 \qquad \E_{v,w\sim\mu}d(v,w)\ge cm.
\end{equation}
The change in either expectation is at most $2C/m$.

Make $k_v$ copies of each vertex. Write $d_0$ for the pulled-back
pseudometric and let $d_{\min}>0$ be the least distance between distinct
original labels. Choose
$0<\zeta\le\min(1,d_{\min})/m^2$ and define
\[
 d'(x,y)=d_0(x,y)+\zeta\mathbf1_{x\ne y}.
\]
The added equilateral metric is squared Hilbert distance, for example
using $\sqrt{\zeta/2}$ times distinct standard unit vectors. Therefore
$d'$ is a genuine negative-type metric, and its uniform metric average
is at least the $\mu$-average in \eqref{eq:mu}.

Here is the full reduction for a map $F$ on the copies with values in
$S_1(H)$, in an order that makes the quantifiers explicit. First collect
the ranges of the finitely many $F(x)$ and $F(x)^*$ into a separable
closed subspace $K\subseteq H$. This is possible because trace-class
operators are compact; their ranges are separable. Moreover
$F(x)=P_KF(x)P_K$. Choose finite-rank orthogonal projections $E_\ell$
increasing strongly to $P_K$ and put $F_\ell(x)=E_\ell F(x)E_\ell$.
Compression is trace-norm contractive, so every $F_\ell$ is nonexpanding
into a finite matrix algebra.

Choose one copy $x_v$ of each label and put $G_\ell(v)=F_\ell(x_v)$.
For distinct labels,
\[
 d'(x_v,x_w)=d(v,w)+\zeta\le(1+m^{-2})d(v,w).
\]
Thus $G_\ell/(1+m^{-2})$ is nonexpanding on the label set. Hermitian
dilation reduces this finite matrix map to the selfadjoint contraction
already proved, so \eqref{eq:mu} applies. Each copy is at distance
$\zeta$ from its representative. It follows that
\begin{equation}\label{eq:copies}
 \E_{x,y\text{ uniform}}\norm{F_\ell(x)-F_\ell(y)}_1
 \le(1+m^{-2})C\sqrt m(\log m)^{7/2}+2\zeta.
\end{equation}
Finally $E_\ell A E_\ell\to A$ in trace norm for every trace-class
$A$ supported on $K$. To see this without a uniform operator-norm
convergence assertion, approximate $A$ in trace norm by a finite-rank
operator. On its finitely many rank-one terms strong convergence of the
projections gives trace-norm convergence; compression contracts the
remaining trace-norm error. Passing to the limit in the finite pair sum
in \eqref{eq:copies} establishes it for $F$. This argument also covers
nonseparable $H$. The reduction order was compression, representatives,
dilation, the finite estimate, and then the trace-norm limit.

By \eqref{eq:size},
\[
 m^m\le n\le\exp(Cm(\log m)^2).
\]
For sufficiently large $m$, this implies $\log\log n\ge\log m$ and
$\sqrt{\log n}\le C\sqrt m\log m$. The ratio of the metric average to
the image average is therefore at least
\[
 c\frac{\sqrt m}{(\log m)^{7/2}}
 \ge c'\frac{\sqrt{\log n}}{(\log\log n)^{9/2}}.
\]
Taking one construction for each sufficiently large $m$ proves
Theorem~\ref{thm:main}. All metric choices precede the choice of $H$ or $F$.

\section{Comparison with other trace-class operations}\label{sec:comparisons}
The precise claim is new to this project, with no antecedent found in a
bounded search. The individual tools above are standard; no global
priority claim follows from the search or from the shortness of the
matrix bridge. We record the operations that distinguish the statement,
rather than claiming that alternative methods cannot ever apply.

\subsection{Dimension reduction, snowflakes and quantum expansion}
Naor--Pisier--Schechtman \cite{NPS} study dimension reduction for subsets
already embedded in trace class, including embeddings into arbitrary
low-dimensional subspaces. Their discussion explains why scalar
coordinate sampling does not automatically have a matricial version.
It does not provide the all-target-dimensions obstruction for our
negative-type source family. Naor--Young \cite{NY}, Question~13 and its
discussion, concern Heisenberg groups and trace-class nonembedding. Our
metrics are different; no conclusion about those groups is proved here.

Kush--Nikolov--Tang \cite{KNT}, Theorem~34 of the cited preprint, embed
the $p/2$-snowflake of Schatten $p$ into Hilbert space with constant
average distortion for $1\le p\le2$. At $p=1$ the source distance has
been replaced by its square root. Composing with a trace-class map of
our negative-type metric therefore controls a Hilbert map for $\sqrt d$,
which is already Euclidean. This does not supply the needed contraction
of first averages for $d$ itself.

Escudero Guti\'errez--Muguruza \cite{EGM} compare expansion notions for
compressed blocks of bistochastic matrix tuples. Their Section~3.2
distinguishes that operation from metric Poincar\'e inequalities for
arbitrary maps. We do not derive an arbitrary-map inequality from channel
expansion, and we do not claim to settle their general Question~2.

\subsection{Talagrand type and nonlinear cube quotients}
Luo--Wang--Zhou \cite{LWW}, Theorem~1.1, give a noncommutative
$L_1$ cube inequality with a sum of coordinate Orlicz increments. Its
input is different from the simultaneous-chart representation controlling
$\E\norm{\sum g_jA_j}_1$ in Lemma~\ref{lem:selector}. We do not claim
that no dimension-free trace-class functional inequalities were known.

Cordero-Erausquin--Eskenazis \cite{CEE}, Theorem~13, combined with
\cite{LWW}, gives a genuine alternative route through quotients of the
cube with small geometric influences. Although their displayed theorem
is a bi-Lipschitz lower bound, equation~(123) in its proof is already a
pair-average inequality, before any co-Lipschitz assumption is inserted
in (124). For a nonexpanding quotient map, the coordinate increment is
bounded by one half of the corresponding boundary indicator, as in
(125). Hence their proof bounds the pair average under the pushforward
of uniform cube measure. It would be incorrect to dismiss this route
because the theorem is phrased in bi-Lipschitz terms.

For precision, write $\mathfrak t$ for the Talagrand-type constant
(to distinguish it from our mesh $\tau$), $\sigma_i$ for the coordinate
boundary probabilities, and $a_1$ for the quotient metric average. For
$\psi(t)=t/\log(e+t)$, the resulting lower bound is
\[
 \frac{2a_1}{\mathfrak t\sum_{i:\sigma_i>0}
                      1/\psi^{-1}(1/\sigma_i)}.
\]
A random coordinate path gives $a_1\le\frac12\sum_i\sigma_i$: each
coordinate is flipped with probability $1/2$, and the intermediate
point is uniform. Also $\psi^{-1}(u)\le2u\log(e+u)$, since
$e+2u\log(e+u)\le(e+u)^2$. Thus the displayed lower bound is at most
\begin{equation}\label{eq:influencecap}
 2\mathfrak t^{-1}\max_{i:\sigma_i>0}\log(e+1/\sigma_i).
\end{equation}
The trivial one-class quotient can be discarded. This is an
influence-logarithmic cap, not an $O(\log\log\mathcal N)$ cap in the
number $\mathcal N$ of quotient classes; no relation forcing the latter
is available. Theorem~42 of \cite{CEE} supplies the same Talagrand-type
property for $L_1$ as \cite{LWW} supplies for $S_1$, with universal
constants (a common upper bound can be used for both).

On our source representation the direct coordinate-sum estimate is
\[
 \sum_{j=1}^m\E\norm{\partial_jH_1}_1
 \le C m\,N m^{-1/2}\frac mN(\log m)^3
 =C m^{3/2}(\log m)^3,
\]
whereas the proof needs $C\sqrt m(\log m)^{7/2}$. The missing gain is
polynomial Gaussian cancellation, not merely a change of constants.
An influence-logarithmic route could supply such a gain only after an
appropriate very-small-influence representation is constructed. The
compared inputs supply no genuine negative-type quotient family with
that influence control, the required average scale, and uniformization
when classes have unequal sizes. They do not identify the multichart
metric with such a quotient. This is a missing construction in the read
inputs, not an impossibility theorem about Talagrand-type arguments.

\subsection{Sharp metric cotype and the usual grids}
Jiao--Luo--Wu \cite{JLW}, posted October~15, 2026, before this
candidate, prove sharp metric cotype $2$ for $S_1$ (Theorem~1.3). In their notation, for every even $k\ge4$ and
$f:\mathbb Z_k^d\to S_1$,
\[
 \sum_{j=1}^d\E_x\norm{f(x+(k/2)e_j)-f(x)}_1^2
 \le(3\sqrt{70}\,k+6\sqrt d)^2
       \E_{x,\epsilon}\norm{f(x+\epsilon)-f(x)}_1^2,
\]
where $\epsilon$ is uniform on $\{-1,0,1\}^d$. Thus the scale can be
$k\asymp\sqrt d$. Their paper also constructs nonlinear matrix
smoothing; its existence is fully credited.

The general-metric cardinality bound below is classical. As recorded in
\cite{GMN}, Section~1.3, Bourgain's 1987 grid theorem gives distortion
$\gtrsim_Y d^{1/q}$ for embeddings of $[\lceil d^{1+1/q}\rceil]^d_\infty$
into a Banach space $Y$ of Rademacher cotype $q$; GMN's Theorem~1.2
also yields this bound. The same section credits Tomczak-Jaegermann
(1974) for cotype $2$ of $S_1$. Here $[m]^d_\infty$ denotes
$\{0,\ldots,m\}^d$ with the sup metric. Thus, for $q=2$,
$\mathcal N=(\lceil d^{3/2}\rceil+1)^d$ and
$\log\mathcal N=(3/2+o(1))d\log d$ give distortion at least
$c\sqrt{\log\mathcal N/\log\log\mathcal N}$.
These two classical attributions are through GMN's explicit statements;
the original Bourgain and Tomczak-Jaegermann bodies were not read here.

The JLW/MN chain recovers this cardinality order with grid side
$O(\sqrt d)$, instead of order $d^{3/2}$. The grid nonembedding deduction
inserts a co-Lipschitz bound for long coordinate jumps into the cotype
inequality; see \cite{MN}, proof of Theorem~1.12. JLW's
Corollary~6.1 has constant $32$ at even scale
$k\ge\max(4,\sqrt d)$. For a nonexpanding map of the sup-metric torus
$\mathbb Z_k^d$ with inverse Lipschitz constant $D$, its left side is at
least $d(k/(2D))^2$, while its right side is at most $(32k)^2$.
Consequently $D\ge\sqrt d/64$. This is the proof of
\cite{MN}, Theorem~1.12, with the new sharp inequality substituted.
The third assertion of \cite{MN}, Lemma~6.12, embeds the torus with
constant distortion into a sup-grid of side $O(k)$ and dimension $O(d)$.
Taking $k\asymp\sqrt d$ gives such grids with $\mathcal N$ points and
trace-class bi-Lipschitz distortion at least
\[
 c\sqrt{\frac{\log\mathcal N}{\log\log\mathcal N}}.
\]
This is the same order in $\mathcal N$ as the classical bound, not an
improvement of that order. The sharp scale improves the grid side,
addressing the $S_1$ case of \cite{GMN}, Question~3. The classical bound
is already stronger than our bi-Lipschitz corollary for general metrics. Any possible novelty of that
corollary is therefore confined to the negative-type specialization,
under the bounded-search qualification above. The grid consequence is
not the first-average negative-type statement of
Theorem~\ref{thm:main}, as the following exact checks show.

\begin{proposition}\label{prop:grid}
For $G=\{0,\ldots,k\}^d$ with $k,d\ge1$ and the sup metric,
there is a nonexpanding map to $S_1(\C)$ preserving at least one third
of the uniform first average. If $d\ge3$ and $k\ge2$, this grid metric
is not of negative type.
\end{proposition}
\begin{proof}
More generally choose a rank-one orthogonal projection $P$ in any
nonzero Hilbert space and set $F(x)=x_1P$. Then
$\norm{F(x)-F(y)}_1=|x_1-y_1|\le\norm{x-y}_\infty$. For independent
uniform $U,V\in\{0,\ldots,k\}$,
\[
 \E|U-V|=\frac2{(k+1)^2}\sum_{h=1}^k h(k+1-h)
       =\frac{k(k+2)}{3(k+1)}\ge k/3.
\]
Since $\E\norm{x-y}_\infty\le k$, the first assertion follows.

For the second assertion take, in the side-two three-dimensional grid,
\[
 A=\{(2,1,0),(0,1,0),(0,1,1)\},\qquad
 B=\{(1,0,1),(1,0,0),(1,2,0)\}.
\]
Assign coefficients $+1$ to $A$ and $-1$ to $B$. Their sum is zero.
Each unordered within-set distance sum is $5$, and all nine cross
distances equal $1$. Thus
\[
 \sum_{i,j}c_ic_j\norm{x_i-x_j}_\infty=2\cdot5+2\cdot5-2\cdot9=2>0.
\]
If these distances were $|z_i-z_j|^2$, the same expression would equal
$-2|\sum_i c_i z_i|^2\le0$. This contradiction persists in every larger
such grid by inclusion.
\end{proof}
The proposition excludes the standard grid substitution, not every use
of metric cotype. To use the cotype inequality for our uniform first
average, one would need a torus map into the negative-type family whose
single-coordinate antipodal pair law controls a fixed fraction of its
uniform pair law, with sufficiently short local jumps. No such map is
supplied by the compared inputs. Cotype itself is a lower-bound property
of random sums, whereas Lemma~\ref{lem:selector} gives an upper bound on
$\E\norm{\sum_jg_jA_j}_1$.

There is also a concrete loss in the nonlinear smoothing alternative.
Proposition~5.2 of \cite{JLW} bounds the $\ell_2$ sum of trace-norm
partial derivatives by $6\sqrt{70}$ times the input oscillation. In
$S_1$, converting that quantity to a Gaussian average can lose
$\sqrt m$: for $A_j=e_je_j^*$,
\[
 \Big(\sum_j\norm{A_j}_1^2\Big)^{1/2}=\sqrt m,
 \qquad \E\norm{\sum_jg_jA_j}_1=m\sqrt{2/\pi}.
\]
That is precisely the size of gain at issue. Moreover, the macro bound
at smoothing scale $\lambda$ gives input oscillation $O(m\lambda)$;
a derivative estimate after rescaling and this $\ell_2$-to-Gaussian
conversion again permits $O(m^{3/2})$. The nonlinear smoothing controls
oscillation, not the simultaneous representation
$A_{sj}=\sum_i u_{sij}b_{si}$ used by the positive spectral selection.
These comparisons do not rule out future combinations of the methods.
The Bessel step in \cite{JLW}, equation~(4.5), and its Hermitian dilation
and finite-rank approximation (Steps~II--III) are parallel standard
devices; they are not claimed as new here.

\section{Scope of the accompanying formal support}\label{sec:formal}
The accompanying author-side Lean development uses Lean~4.34.1 and
mathlib revision
\begin{center}\small\texttt{d13f23b723b8a846827a245b89c10fc7d3f11612}.\end{center}
It contains 31 support declarations and 15 control declarations. For
arbitrary finite complex positive matrices $M_s$ it constructs actual
spectral projections and proves their defect and compression bounds,
the half-coverage conclusion from the summed budget, and the positive
trace pairing without a commutation assumption. It also proves the
per-chart budget consequence. The controls include noncommuting
rank-one projections, necessity of positivity in the trace pairing,
a zero-threshold failure for the abstract defect interface, and exact
threshold endpoint behavior. These are general kernel-checked lemmas
and hypothesis tests, not a numerical proof of the metric theorem.

The support does not formalize the Wick/register argument, Gaussian
integration or finite direction construction; fractional powers and the
Loewner--Heinz/concavity input; the conversion of compressed squares to
operator norms; the Gaussian construction of $Y$ or trace-norm duality;
source geometry, coarea, copy uniformization, asymptotics or arbitrary-$H$
compression. Its positive matrix and summed-budget hypotheses still
have to be supplied by the mathematical proof above. The build audits
all declaration types and allows only the standard axioms
\texttt{propext}, \texttt{Classical.choice} and \texttt{Quot.sound}.
The pinned dependency object cache is trusted, not rebuilt. A separate
independent review passed this exact scoped support, with no required
repairs. Its additional checks exercise actual coverage and trace
interfaces on noncommuting data. In particular, the genuine zero-matrix
spectral cutoff at threshold zero is $I$; the abstract zero-threshold
counterexample is not a counterexample to that cutoff. The chosen
eigenbasis is used to construct the projection; independence of that
choice or identification with continuous functional calculus is not
formalized. Only the proved projection properties are used here. Neither
the 46 author declarations nor the review constitute a formal proof of
Theorem~\ref{thm:main}.

\paragraph{Reproduction.}
The accompanying packet README gives the PDF command
\texttt{python3 -B build.py --output /var/tmp/NEW-PDF-BUILD}, with a
fresh output directory and an explicit Tectonic option if needed.
The separate scoped formal packet \path{t117q-formal-v1/rebuild.py}
provides the Lean entry point with \texttt{--output}, \texttt{--lean}
and \texttt{--cache} arguments; its README specifies the pinned tools
and trusted dependency cache. These builds do not replace the stated
mathematical and formal-coverage limits.

\appendix
\section{Exact source interfaces and reading limits}\label{sec:source}
The source \cite{source} is pinned at OpenAI/math revision
\begin{center}\small\path{fd4aeeb2ee4fc729c18d98444fed42fd0529eeeb}.\end{center}
Its path is
\begin{center}\small\path{preprints/Near-square-root-logarithmic-integrality-gaps-for-uniform-sparsest-cut-September-24-2026/build/source/}.\end{center}
The complete preprint-tree inventory, not a truncated root-tree listing,
identifies the inputs. Source truth is stipulated. The exact reused
interfaces are:
\begin{itemize}
\item \path{directions.tex}: \path{lem:directions}, its norm,
covariance and Fourier approximation parts, and
\path{eq:pseudoinverse}. The original $S=m^3$ assertion is not used
unchanged; Section~\ref{sec:design} rechecks the probability bounds in
\path{directions-proof.tex}. Scalar majority charts are unused.
\item \path{kernel.tex}: \path{eq:kernel},
\path{ker:A-definition}, \path{ker:extension-identity},
\path{ker:profile}, \path{ker:strong-bounds},
\path{ker:P-definition} and \path{ker:DP-identity}.
Section~\ref{sec:geometry} reproduces the formulas and checks the new
$S,\varepsilon,\tau$ dependence explicitly.
\item \path{repair.tex}: \path{lem:repair}, including
\path{rep:eq:hypothesis}, \path{rep:eq:distance} and
\path{rep:eq:local}, stated for an arbitrary finite chart partition.
\item \path{metric.tex}: the counting, diameter and average arguments
behind \path{met:eq:macro}, \path{met:eq:adjacent},
\path{met:eq:size-diameter} and \path{met:eq:average}.
Their modified bounds are proved here, not imported at the old mesh.
\item \path{contraction.tex}: smoothing, derivative mismatch,
distributional gradient, coarea coefficient formula, density variation
and lattice-sum bound. Its scalar-coordinate good-chart step is replaced
by Lemma~\ref{lem:selector}; its Gaussian rotation is used in the
norm-valued form proved in Lemma~\ref{lem:cube}.
\item \path{reduction.tex}: the standard finite replication idea.
Section~\ref{sec:completion} proves the precise uniformization used here
and adds an equilateral term to retain strict positivity.
\end{itemize}
The constant-factor-hardness companion was keyword-screened, not read
through its operative proof, and is not a dependency. No assertion about
that full proof is used. Tropp's matrix-concentration text, accessed only
in earlier browser passages without an authenticated local body, is
orientation rather than evidence and is not needed here. The classical
operator Khintchine and design references are credits to established
methods, not assertions that their complete proofs were read in this
manuscript task. The new lemmas needed here have their proofs above.

Primary comparisons are limited to the stated theorems and operative
passages of \cite{NPS,NY,KNT,LWW,CEE,JLW,MN,GMN,EGM}; no exhaustive
bibliography or forward-citation search was performed. The full
mathematical proof passed a separate proof review, and contribution
round~1 classified the new bridge as substantial to the project while
requiring the positioning implemented in Section~\ref{sec:comparisons}.
A separate round~2 contribution review retained that project-level
classification, with the corrected influence cap, concrete smoothing
loss and stronger general-metric grid consequence now recorded in
Section~\ref{sec:comparisons}. Neither review establishes literature
priority or constitutes manuscript acceptance.

\paragraph{Acknowledgments and license.}
Parent author collaboration checked the register-contraction and selector
interfaces and supplied the explicit grid-average and six-point
certificates in Proposition~\ref{prop:grid}. This assistance is author
work, not independent review. Original prose is licensed CC BY~4.0 and
original accompanying code MIT. Third-party manuscripts and their
licenses are separate; private comparison bodies are not redistributed.

\begin{thebibliography}{99}
\bibitem{source} OpenAI, \emph{Near-square-root logarithmic integrality gaps
for uniform sparsest cut}, September 24, 2026, pinned manuscript.
\url{https://github.com/openai/math/tree/fd4aeeb2ee4fc729c18d98444fed42fd0529eeeb/preprints/Near-square-root-logarithmic-integrality-gaps-for-uniform-sparsest-cut-September-24-2026}.
\bibitem{HP} U. Haagerup and G. Pisier, \emph{Bounded linear operators
between $C^*$-algebras}, 1993. \url{https://arxiv.org/abs/math/9302214}.
\bibitem{Buchholz} A. Buchholz, \emph{Operator Khintchine inequality in
non-commutative probability}, Math. Ann. \textbf{319} (2001), 1--16.
\url{https://doi.org/10.1007/PL00004425}.
\bibitem{SZ} P. D. Seymour and T. Zaslavsky, \emph{Averaging sets:
a generalization of mean values and spherical designs}, Adv. Math.
\textbf{52} (1984), 213--240.
\url{https://doi.org/10.1016/0001-8708(84)90022-7}.
\bibitem{Pisier} G. Pisier, \emph{Probabilistic methods in the geometry
of Banach spaces}, in Probability and Analysis, Lecture Notes in Math.
\textbf{1206}, Springer, 1986, 167--241.
\url{https://doi.org/10.1007/BFb0076302}.
\bibitem{NPS} A. Naor, G. Pisier and G. Schechtman, \emph{Impossibility
of dimension reduction in the nuclear norm}, preprint, 2017;
conference version SODA 2018. \url{https://arxiv.org/abs/1710.08896v1}.
\bibitem{NY} A. Naor and R. Young, \emph{Vertical perimeter versus
horizontal perimeter}, Ann. of Math. \textbf{188} (2018), 171--279.
\url{https://arxiv.org/abs/1701.00620}.
\bibitem{KNT} D. Kush, A. Nikolov and H. Tang, \emph{Near neighbor
search via efficient average distortion embeddings}, SoCG 2021;
full preprint \url{https://arxiv.org/abs/2105.04712v1}.
\bibitem{LWW} S. Luo, Y. Wang and D. Zhou, \emph{Talagrand type for
noncommutative $L_1$ spaces}, 2026.
\url{https://arxiv.org/abs/2610.03407v1}.
\bibitem{CEE} D. Cordero-Erausquin and A. Eskenazis,
\emph{Talagrand's influence inequality revisited}, cited preprint version,
2023. \url{https://arxiv.org/abs/2010.11583v3}.
\bibitem{JLW} Y. Jiao, S. Luo and L. Wu, \emph{The sharp scale inequality
for the Schatten trace class}, October 15, 2026.
\url{https://arxiv.org/abs/2610.21445v1}.
\bibitem{MN} M. Mendel and A. Naor, \emph{Metric cotype}, Ann. of Math.
\textbf{168} (2008), 247--298; cited preprint version
\url{https://arxiv.org/abs/math/0506201v5}.
\bibitem{GMN} O. Giladi, M. Mendel and A. Naor, \emph{Improved bounds
in the metric cotype inequality for Banach spaces}, cited preprint version,
2010. \url{https://arxiv.org/abs/1003.0279v3}.
Section~1.3 and its references [3] and [18] are the attribution source
used here for Bourgain's 1987 grid theorem and Tomczak-Jaegermann's
1974 cotype theorem; their original bodies were not consulted.
\bibitem{EGM} F. Escudero Guti\'errez and G. Muguruza,
\emph{All $S_p$ notions of quantum expansion are equivalent}, 2024.
\url{https://arxiv.org/abs/2405.03517v1}.
\end{thebibliography}
\end{document}
